\documentclass[12pt, a4paper]{article}

% ── Encoding & language ────────────────────────────────────────────────────
\usepackage[utf8]{inputenc}
\usepackage[T1]{fontenc}
\usepackage[english]{babel}

% ── Mathematics ────────────────────────────────────────────────────────────
\usepackage{amsmath, amssymb, amsthm}
\usepackage{mathtools}          % dcases, prescript, etc.
\usepackage{physics}            % \bra, \ket, \norm, \abs, \dv, \pdv

% ── Figures & layout ───────────────────────────────────────────────────────
\usepackage{graphicx}
\usepackage{float}
\usepackage{subcaption}
\usepackage[margin=2.5cm]{geometry}
\usepackage{setspace}
\onehalfspacing

% ── Hyperlinks ─────────────────────────────────────────────────────────────
\usepackage[colorlinks=true, linkcolor=blue!70!black,
            citecolor=blue!70!black, urlcolor=blue!70!black]{hyperref}

% ── Code listings ──────────────────────────────────────────────────────────
\usepackage{listings}
\usepackage{xcolor}
\lstset{
  basicstyle=\ttfamily\small,
  keywordstyle=\color{blue!80!black}\bfseries,
  commentstyle=\color{gray},
  stringstyle=\color{green!50!black},
  breaklines=true,
  frame=single,
  numbers=left,
  numberstyle=\tiny\color{gray}
}

% ── Custom commands ────────────────────────────────────────────────────────
\newcommand{\dt}{\Delta t}
\newcommand{\dx}{\Delta x}
\newcommand{\hb}{\hbar}
\newcommand{\ii}{\mathrm{i}}          % imaginary unit (upright)
\newcommand{\ee}{\mathrm{e}}          % Euler's number (upright)
\newcommand{\CN}{Crank--Nicolson}
\newcommand{\Htop}{\hat{H}}
\newcommand{\Ktop}{\hat{K}}
\newcommand{\Vtop}{\hat{V}}

% ── Bibliography ───────────────────────────────────────────────────────────
\usepackage[numbers, sort&compress]{natbib}

% ───────────────────────────────────────────────────────────────────────────
\begin{document}

% ── Title page ─────────────────────────────────────────────────────────────
\begin{titlepage}
  \centering
  \vspace*{3cm}
  {\LARGE\bfseries Numerical Solution of the Time-Dependent\\[6pt]
   Schr\"{o}dinger Equation\\[4pt]
   via the Crank--Nicolson Scheme\par}
  \vspace{1.5cm}
  {\large Quantum tunnelling of a Gaussian wave packet\\
   through a rectangular potential barrier\par}
  \vspace{2cm}
  {\large Ludovico Fecia di Cossato\par}
  \vfill
  {\large\itshape TDSE\_CN — theoretical notes\par}
  \vspace{0.5cm}
  {\normalsize Atomic units throughout: $\hbar = m_e = e = 1$\par}
  \vspace{1cm}
  {\normalsize \today\par}
\end{titlepage}

\tableofcontents
\newpage

% ═══════════════════════════════════════════════════════════════════════════
\section{Introduction}
% ═══════════════════════════════════════════════════════════════════════════

The time-dependent Schr\"{o}dinger equation (TDSE) governs the evolution of
any non-relativistic quantum system.  Even in one spatial dimension, it
exhibits phenomena — wave packet spreading, interference, and quantum
tunnelling — that have no classical analogue and are central to modern
physics and technology (tunnel diodes, scanning tunnelling microscopy,
nuclear alpha decay, \ldots).

Despite its conceptual depth, the one-dimensional TDSE with a static
potential is fully tractable numerically: the problem reduces to propagating
a complex vector in time via a sequence of sparse linear systems.  This
makes it an ideal testbed for implicit finite-difference methods and for
developing intuition about quantum dynamics.

These notes describe the physical setup, the mathematical structure of the
\CN{} discretisation, and the Thomas algorithm used to solve the resulting
tridiagonal system at each time step.  A final section analyses the output
of the numerical simulation qualitatively and quantitatively.

% ═══════════════════════════════════════════════════════════════════════════
\section{Physical setup}
% ═══════════════════════════════════════════════════════════════════════════

\subsection{The time-dependent Schr\"{o}dinger equation}

The state of a one-dimensional quantum particle is fully described by a
complex-valued wave function $\psi(x,t) \in L^2(\mathbb{R})$.  Its
time evolution is governed by the TDSE
\begin{equation}
  \ii\hbar\,\pdv{\psi}{t}(x,t) \;=\; \Htop\,\psi(x,t),
  \label{eq:tdse}
\end{equation}
where the Hamiltonian operator $\Htop = \Ktop + \Vtop$ splits into a
kinetic part
\begin{equation}
  \Ktop \;=\; -\frac{\hbar^2}{2m}\pdv[2]{}{x}
\end{equation}
and a potential energy operator $\Vtop = V(x)$ acting multiplicatively.

In \textbf{atomic units} ($\hbar = m_e = 1$) — used throughout this
document — equation \eqref{eq:tdse} becomes
\begin{equation}
  \ii\,\pdv{\psi}{t} \;=\;
  \left[-\frac{1}{2}\pdv[2]{}{x} + V(x)\right]\psi.
  \label{eq:tdse_au}
\end{equation}

\subsection{Physical interpretation of the wave function}

The wave function $\psi(x,t)$ is not directly observable.  The physically
meaningful quantity is the \emph{probability density}
\begin{equation}
  \rho(x,t) \;=\; |\psi(x,t)|^2,
\end{equation}
which gives the probability of finding the particle between $x$ and $x+dx$
as $\rho(x,t)\,dx$.  The total probability is conserved:
\begin{equation}
  \int_{-\infty}^{+\infty} |\psi(x,t)|^2\,dx \;=\; 1
  \quad \forall\, t.
  \label{eq:norm}
\end{equation}
This constraint — the \emph{norm conservation} or \emph{unitarity} of the
evolution — is one of the most important consistency checks for any
numerical scheme.

\subsection{The Gaussian wave packet}

The initial state chosen for this simulation is a \emph{Gaussian wave
packet}, the quantum state that most closely mimics a classical particle
with a definite position and momentum:
\begin{equation}
  \psi(x,0) \;=\; A\,\exp\!\left[-\frac{(x-x_0)^2}{4\sigma^2}\right]
  \exp\!\left(\ii k_0 x\right),
  \label{eq:gauss}
\end{equation}
where:
\begin{itemize}
  \item $x_0$ is the initial position of the packet centre,
  \item $\sigma$ is the spatial width (standard deviation of $\rho$),
  \item $k_0$ is the central wave vector, corresponding to mean momentum
        $\langle p\rangle = \hbar k_0$,
  \item $A = (2\pi\sigma^2)^{-1/4}$ is the normalisation constant.
\end{itemize}

The Gaussian is the \emph{minimum-uncertainty state}: it saturates the
Heisenberg uncertainty relation $\Delta x\,\Delta p \geq \hbar/2$.  Its
Fourier transform is also a Gaussian, centred at $k_0$ with width
$1/(2\sigma)$, so the momentum spread is $\Delta p = \hbar/(2\sigma)$.

The mean kinetic energy of the packet is
\begin{equation}
  E \;=\; \frac{\langle p\rangle^2}{2m} \;=\; \frac{k_0^2}{2}
  \quad \text{(atomic units)}.
  \label{eq:energy}
\end{equation}
This is the quantity to compare against the barrier height $V_0$ in order
to classify the dynamical regime.

\subsection{The rectangular potential barrier}

The external potential is a rectangular barrier of height $V_0$ and width
$w$ located at $x \in [x_\text{bar},\, x_\text{bar}+w]$:
\begin{equation}
  V(x) \;=\;
  \begin{cases}
    V_0 & \text{if } x_\text{bar} < x < x_\text{bar} + w, \\
    0   & \text{otherwise.}
  \end{cases}
  \label{eq:barrier}
\end{equation}
The barrier is the simplest non-trivial scattering potential and admits an
exact analytic solution for plane-wave states (see
Section~\ref{sec:tunnelling}), making it ideal for validating the
numerics.

% ═══════════════════════════════════════════════════════════════════════════
\section{Quantum tunnelling}
\label{sec:tunnelling}
% ═══════════════════════════════════════════════════════════════════════════

\subsection{Stationary analysis}

For a monochromatic plane wave $\psi(x) = \ee^{\ii kx}$ with energy
$E = k^2/2$, the Schr\"{o}dinger equation can be solved piecewise in
each region.  Matching wave functions and their derivatives at the
barrier edges yields the exact transmission coefficient \cite{Griffiths2018}
\begin{equation}
  T \;=\; \left[1 + \frac{V_0^2\sinh^2(\kappa w)}{4E(V_0-E)}\right]^{-1}
  \qquad (E < V_0),
  \label{eq:T_tunnelling}
\end{equation}
where $\kappa = \sqrt{2(V_0-E)}$ (atomic units) is the evanescent decay
constant inside the barrier.  The reflection coefficient is simply
$R = 1 - T$.

Several qualitative features are immediately apparent from
\eqref{eq:T_tunnelling}:
\begin{enumerate}
  \item $T > 0$ even when $E < V_0$ — this is \emph{quantum tunnelling}.
        Classically, the particle cannot penetrate a region where $E < V$.
  \item $T$ decays \emph{exponentially} with barrier width: for
        $\kappa w \gg 1$,
        \begin{equation}
          T \;\approx\; 16\,\frac{E}{V_0}\!\left(1 - \frac{E}{V_0}\right)
          \ee^{-2\kappa w}.
        \end{equation}
  \item In the above-barrier regime $E > V_0$, partial reflection still
        occurs due to the wave nature of $\psi$, with
        \begin{equation}
          T \;=\; \left[1 + \frac{V_0^2\sin^2(k' w)}{4E(E-V_0)}\right]^{-1},
          \quad k' = \sqrt{2(E-V_0)}.
        \end{equation}
        Resonant transmission ($T = 1$) occurs whenever $k'w = n\pi$.
\end{enumerate}

\subsection{Wave-packet dynamics}

A Gaussian wave packet is a superposition of plane waves:
\begin{equation}
  \psi(x,0) \;=\; \int_{-\infty}^{+\infty}
  \tilde{\phi}(k)\,\ee^{\ii kx}\,\frac{dk}{2\pi},
  \qquad
  \tilde{\phi}(k) \propto \ee^{-(k-k_0)^2\sigma^2}\,\ee^{-\ii k x_0}.
\end{equation}
Each component $k$ evolves independently and interacts with the barrier
with its own coefficient $T(k)$.  After the scattering event, the total
wave function splits into a \emph{reflected} part travelling to the left
and a \emph{transmitted} part travelling to the right.

The integrated probabilities
\begin{align}
  R(t) &= \int_{-\infty}^{x_\text{bar}} |\psi(x,t)|^2\,dx, \\
  T(t) &= \int_{x_\text{bar}}^{+\infty} |\psi(x,t)|^2\,dx
\end{align}
approach constant asymptotic values after the scattering is complete.
These asymptotic values can be compared to the plane-wave formula
\eqref{eq:T_tunnelling} evaluated at $k = k_0$ — agreement is expected
when the packet is narrow in $k$-space (large $\sigma$).

% ═══════════════════════════════════════════════════════════════════════════
\section{Numerical method}
% ═══════════════════════════════════════════════════════════════════════════

\subsection{Spatial discretisation}

The domain $[x_\text{min},\, x_\text{max}]$ is divided into $N$ uniformly
spaced points with grid spacing $\dx = L/(N-1)$, where $L =
x_\text{max} - x_\text{min}$.  The wave function is represented as the
vector $\boldsymbol{\psi}^n \in \mathbb{C}^N$ with components
$\psi_i^n \approx \psi(x_i, n\dt)$.

The second derivative is replaced by the standard symmetric
second-order finite-difference stencil:
\begin{equation}
  \pdv[2]{\psi}{x}\bigg|_{x_i}
  \;\approx\;
  \frac{\psi_{i-1} - 2\psi_i + \psi_{i+1}}{\dx^2}
  + \mathcal{O}(\dx^2).
  \label{eq:fd_stencil}
\end{equation}
This approximates the kinetic part of $\Htop$ as a real symmetric
tridiagonal matrix $K$ with main diagonal $-2/\dx^2$ and off-diagonals
$+1/\dx^2$ (multiplied by $-1/2$):
\begin{equation}
  [H_\text{disc}]_{ij}
  \;=\;
  -\frac{1}{2\dx^2}\begin{pmatrix}
    -2 & 1 \\
    1 & -2 & 1 \\
      & \ddots & \ddots & \ddots \\
      &        & 1 & -2 & 1 \\
      &        &   & 1  & -2
  \end{pmatrix}_{ij}
  + V_i\,\delta_{ij}.
  \label{eq:Hdisc}
\end{equation}

\subsection{Time discretisation: the Crank--Nicolson scheme}

A naive explicit (forward-Euler) discretisation of the TDSE,
\begin{equation}
  \boldsymbol{\psi}^{n+1} = \boldsymbol{\psi}^n
  - \ii\dt\, H_\text{disc}\,\boldsymbol{\psi}^n,
\end{equation}
is \emph{conditionally stable}: the time step must satisfy
$\dt \lesssim \dx^2$, which for $\dx = 0.01$\,bohr means
$\dt \lesssim 10^{-4}$\,a.u. — two orders of magnitude smaller than the
\CN{} step used here.  Moreover, the explicit scheme does not conserve
the norm exactly.

The \textbf{\CN{} scheme} \cite{CrankNicolson1947} is an implicit
second-order method that averages the right-hand side between the current
and the next time level:
\begin{equation}
  \ii\,\frac{\boldsymbol{\psi}^{n+1} - \boldsymbol{\psi}^n}{\dt}
  \;=\;
  \frac{H_\text{disc}}{2}\!\left(\boldsymbol{\psi}^{n+1}
  + \boldsymbol{\psi}^n\right).
  \label{eq:cn_raw}
\end{equation}
Rearranging \eqref{eq:cn_raw}:
\begin{equation}
  \underbrace{\left(I + \frac{\ii\dt}{2}\,H_\text{disc}\right)}_{A}
  \boldsymbol{\psi}^{n+1}
  \;=\;
  \underbrace{\left(I - \frac{\ii\dt}{2}\,H_\text{disc}\right)}_{B}
  \boldsymbol{\psi}^n,
  \label{eq:cn_system}
\end{equation}
which is a linear system $A\,\boldsymbol{\psi}^{n+1} = \mathbf{b}^n$
with $\mathbf{b}^n = B\,\boldsymbol{\psi}^n$.

\subsection{Properties of the Crank--Nicolson scheme}
\label{sec:cn_properties}

\paragraph{Unconditional stability.}
The eigenvalues of $H_\text{disc}$ are real (it is Hermitian).  For an
eigenvalue $\lambda$, the amplification factor of the \CN{} scheme is
\begin{equation}
  g \;=\; \frac{1 - \ii\dt\lambda/2}{1 + \ii\dt\lambda/2},
\end{equation}
which satisfies $|g| = 1$ for \emph{any} $\dt > 0$.  The scheme is
therefore \textbf{unconditionally stable}: no time-step restriction is
required for stability.

\paragraph{Norm (probability) conservation.}
The matrices $A$ and $B$ satisfy $B = A^\dagger$ (complex conjugate
transpose).  The discrete propagator $A^{-1}B$ is therefore
\emph{unitary}:
\begin{equation}
  (A^{-1}B)^\dagger (A^{-1}B)
  \;=\; B^\dagger (A^\dagger)^{-1} A^{-1} B
  \;=\; A A^{-1} (A^\dagger)^{-1} B
  \;=\; I.
\end{equation}
Consequently, $\|\boldsymbol{\psi}^{n+1}\|^2 = \|\boldsymbol{\psi}^n\|^2$
at every step — the discrete norm is exactly conserved (up to
floating-point rounding errors of order $10^{-14}$).

\paragraph{Second-order accuracy.}
The local truncation error of \eqref{eq:cn_raw} is $\mathcal{O}(\dt^2)$
in time and $\mathcal{O}(\dx^2)$ in space, so the scheme is globally
second-order in both variables.

\subsection{Matrix structure}

The matrices $A$ and $B$ are \emph{complex tridiagonal}.  Defining the
parameter
\begin{equation}
  r \;=\; \frac{\dt}{4\dx^2},
\end{equation}
their non-zero elements are (for $1 \leq i \leq N-2$):
\begin{alignat}{2}
  A_{i,i}   &= 1 + \ii\frac{\dt}{2}\!\left(\frac{1}{\dx^2} + V_i\right),
  &\quad A_{i,i\pm 1} &= -\ii r, \label{eq:A_elements}\\[4pt]
  B_{i,i}   &= 1 - \ii\frac{\dt}{2}\!\left(\frac{1}{\dx^2} + V_i\right),
  &\quad B_{i,i\pm 1} &= +\ii r. \label{eq:B_elements}
\end{alignat}
The first and last rows enforce Dirichlet boundary conditions
$\psi_0 = \psi_{N-1} = 0$: both rows are set to the identity row
$(1,\,0,\,\ldots)$ in $A$ and $B$.

Because $V(x)$ is static, $A$ and $B$ are \textbf{built once} before
the time loop and reused at every step.

\subsection{The Thomas algorithm}
\label{sec:thomas}

At each time step, the right-hand side is evaluated as a matrix-vector
product $\mathbf{b}^n = B\,\boldsymbol{\psi}^n$ (one pass over the grid,
$\mathcal{O}(N)$ operations), and the resulting tridiagonal system
$A\,\boldsymbol{\psi}^{n+1} = \mathbf{b}^n$ is solved with the
\textbf{Thomas algorithm} \cite{Thomas1949}.

The Thomas algorithm is Gaussian elimination specialised for tridiagonal
systems.  For a system with lower diagonal $l_i$, main diagonal $d_i$,
upper diagonal $u_i$, and right-hand side $b_i$:

\paragraph{Forward sweep} (elimination of $l_i$, from $i=1$ to $N-1$):
\begin{align}
  w_i &= l_i / d_{i-1}, \\
  d_i &\leftarrow d_i - w_i\,u_{i-1}, \\
  b_i &\leftarrow b_i - w_i\,b_{i-1}.
\end{align}

\paragraph{Back substitution} (from $i=N-1$ down to $0$):
\begin{align}
  x_{N-1} &= b_{N-1}/d_{N-1}, \\
  x_i     &= (b_i - u_i\,x_{i+1}) / d_i.
\end{align}

The total cost is $\mathcal{O}(N)$ in both time and memory — optimal for
a tridiagonal system.  The algorithm is numerically stable for
diagonally dominant matrices, which the \CN{} matrix $A$ always satisfies
(the main diagonal dominates because $|A_{ii}| > 2|A_{i,i\pm 1}|$
whenever $\dt/\dx^2 < 2$, easily satisfied in practice).

This is the same algorithm implemented by LAPACK's routine
\texttt{zgtsv}, here realised explicitly in complex arithmetic without
any external library dependency.

\subsection{Boundary conditions}

\textbf{Dirichlet} boundary conditions $\psi(x_\text{min},t) =
\psi(x_\text{max},t) = 0$ are imposed.  Physically, the walls act as
infinite potential barriers.  They introduce artificial reflections if
the wave packet reaches the boundary, so the domain length $L$ must be
chosen large enough that the packet remains far from the walls throughout
the simulation.

A practical criterion is $x_\text{min} \ll x_0 - k_0 T_\text{tot}$ and
$x_\text{max} \gg x_0 + k_0 T_\text{tot}$, where $k_0 T_\text{tot}$ is
the approximate distance travelled by the centre of the packet.

\subsection{Algorithm summary}

The complete time-stepping algorithm is:

\begin{enumerate}
  \item Read parameters and build grid $\{x_i\}$ and potential $\{V_i\}$.
  \item Initialise $\boldsymbol{\psi}^0$ as the normalised Gaussian
        \eqref{eq:gauss}.
  \item Build the tridiagonal matrices $A$ and $B$ (once).
  \item \textbf{For} $n = 0, 1, \ldots, M-1$:
  \begin{enumerate}
    \item[(a)] Compute $\mathbf{b}^n = B\,\boldsymbol{\psi}^n$
               (matrix-vector product, $\mathcal{O}(N)$).
    \item[(b)] Solve $A\,\boldsymbol{\psi}^{n+1} = \mathbf{b}^n$
               with Thomas algorithm ($\mathcal{O}(N)$).
    \item[(c)] Enforce $\psi_0^{n+1} = \psi_{N-1}^{n+1} = 0$.
    \item[(d)] Write snapshot and norm to file (if snapshot step).
  \end{enumerate}
\end{enumerate}
Total cost per time step: $\mathcal{O}(N)$.  Total cost: $\mathcal{O}(MN)$.

% ═══════════════════════════════════════════════════════════════════════════
\section{Analysis of the simulation output}
% ═══════════════════════════════════════════════════════════════════════════

This section describes the expected physical behaviour and the
corresponding signatures in each of the three plots produced by the
Python analysis scripts.

The default simulation parameters are:
\begin{center}
\begin{tabular}{lll}
\hline
Parameter & Value & Role \\
\hline
$L$ & 20\,bohr & domain length \\
$\dx$ & 0.01\,bohr & spatial step \\
$\dt$ & 0.001\,a.u. & time step \\
$T_\text{tot}$ & 2.0\,a.u. & total time \\
$x_0$ & $-5.0$\,bohr & initial packet centre \\
$\sigma$ & 0.5\,bohr & packet width \\
$k_0$ & 5.0\,bohr$^{-1}$ & central wave vector \\
$x_\text{bar}$ & 0.0\,bohr & barrier left edge \\
$w$ & 0.5\,bohr & barrier width \\
$V_0$ & 20.0\,$E_h$ & barrier height \\
\hline
\end{tabular}
\end{center}

The mean kinetic energy of the packet is $E = k_0^2/2 = 12.5\,E_h$, while
the barrier height is $V_0 = 20.0\,E_h$, so $E < V_0$: the simulation is
in the \textbf{tunnelling regime}.  The evanescent decay constant inside
the barrier is $\kappa = \sqrt{2(V_0-E)} \approx 3.87$\,bohr$^{-1}$,
giving an approximate transmission probability from
\eqref{eq:T_tunnelling}:
\begin{equation}
  T_\text{analytic} \;\approx\;
  \left[1 + \frac{20^2\sinh^2(3.87\times 0.5)}{4\times 12.5\times 7.5}
  \right]^{-1}
  \;\approx\; 0.04.
  \label{eq:T_predicted}
\end{equation}
About 4\% of the probability is expected to tunnel through the barrier,
with the remaining 96\% reflected.

\subsection{Snapshot plot (\texttt{snapshots.png})}

\paragraph{What is shown.}
Five panels display $|\psi(x,t)|^2$ at equally-spaced times
$t = 0,\, T/4,\, T/2,\, 3T/4,\, T$.
The potential barrier is shown as a shaded region scaled to the
maximum probability density for visual reference.

\paragraph{Expected behaviour.}
\begin{itemize}
  \item \textbf{$t = 0$}: A single Gaussian peak centred at $x_0 = -5$\,bohr,
        well to the left of the barrier.  Its width is $\sigma = 0.5$\,bohr.

  \item \textbf{Propagation phase ($t \approx T/4$)}: The packet moves
        rightward at group velocity $v_g = k_0 = 5$\,bohr/a.u.  It
        gradually broadens due to dispersion: because different $k$
        components travel at slightly different speeds $v(k) = k$, the
        packet spreads as $\sigma(t) \approx \sqrt{\sigma^2 + t^2/(4\sigma^2)}$.

  \item \textbf{Scattering ($t \approx T/2$)}: The packet reaches the
        barrier.  The probability density is temporarily compressed and
        distorted — the leading edge decelerates while the tail catches up.
        A small transmitted fraction leaks through the right side of the
        barrier.

  \item \textbf{Post-scattering ($t \approx 3T/4$ and $T$)}: Two distinct
        peaks are visible.  A large reflected packet travels leftward; a
        small transmitted packet continues rightward.  The relative areas
        under the two peaks reflect the ratio $R:T \approx 96:4$.
\end{itemize}

\paragraph{What to look for.}
The total area under the curve (norm) must remain constant across all
panels.  Any visible growth or decay signals a numerical problem.

\subsection{Quantitative analysis plot (\texttt{analysis.png})}

\subsubsection{Panel A: Norm conservation}

\paragraph{What is shown.}
The discrete norm $\|\boldsymbol{\psi}^n\|^2_\dx = \dx\sum_i |\psi_i^n|^2$
as a function of time, compared against the exact value of~1.

\paragraph{Expected behaviour.}
Because the \CN{} propagator is unitary (Section~\ref{sec:cn_properties}),
the norm deviation should be $|\|\psi^n\|^2 - 1| \lesssim 10^{-12}$
throughout — essentially machine precision.  This is fundamentally
different from explicit schemes, which dissipate or amplify the norm.

\paragraph{What to look for.}
A flat line at 1 confirms the unitarity of the numerical propagator.
Any monotonic drift would indicate an implementation bug; oscillatory
drift would suggest stability issues.

\subsubsection{Panel B: Transmission and reflection coefficients}

\paragraph{What is shown.}
The time evolution of
\begin{equation}
  R(t) = \dx\sum_{i:\,x_i < x_\text{bar}} |\psi_i|^2,
  \qquad
  T(t) = \dx\sum_{i:\,x_i \geq x_\text{bar}} |\psi_i|^2,
\end{equation}
together with their sum $R(t) + T(t)$.

\paragraph{Expected behaviour.}
Initially $R \approx 1$ and $T \approx 0$ (entire packet to the left of
the barrier).  As the packet reaches and interacts with the barrier, $R$
decreases and $T$ increases.  After the scattering is complete, both
coefficients reach their asymptotic values: $T_\infty \approx 0.04$ and
$R_\infty \approx 0.96$ from \eqref{eq:T_predicted}.

The sum $R(t) + T(t)$ should remain equal to 1 at all times (it is just
another way of writing the norm conservation), providing a self-consistency
check.  During the scattering phase, $R+T$ may appear to dip very slightly
if the wave function is spread \emph{across} the barrier boundary, but
this is a boundary effect of the integration region, not a physical violation.

\paragraph{What to look for.}
The asymptotic value of $T_\infty$ can be directly compared to the
analytic plane-wave result \eqref{eq:T_predicted}.  Discrepancies arise
because the Gaussian packet has a finite momentum spread; the wave-packet
result converges to the plane-wave formula in the limit $\sigma \to \infty$
(narrow $k$-distribution).

\subsubsection{Panel C: Mean position $\langle x \rangle(t)$}

\paragraph{What is shown.}
The expectation value of position,
\begin{equation}
  \langle x\rangle(t) \;=\; \dx\sum_i x_i\,|\psi_i(t)|^2,
\end{equation}
as a function of time.

\paragraph{Expected behaviour.}
\begin{itemize}
  \item \textbf{Before the barrier}: $\langle x\rangle$ increases linearly
        with slope $v_g = k_0 = 5$\,bohr/a.u., reflecting the classical
        motion of the packet centre.

  \item \textbf{During scattering}: the slope decreases as the packet
        interacts with the barrier.  Classically, the particle would be
        reflected instantaneously; quantum mechanically, the interaction is
        spread over a finite time (related to the \emph{Wigner delay time}).

  \item \textbf{After scattering}: $\langle x\rangle$ is the
        probability-weighted average of the reflected and transmitted
        centres.  Since $R \gg T$, the mean position moves leftward (the
        reflected packet dominates).  The slope of the asymptotic linear
        regime will be negative, approximately $-k_0 R_\infty$.
\end{itemize}

\paragraph{What to look for.}
The transition from positive to negative slope signals the moment of
scattering.  The asymptotic slope encodes the relative weight of reflected
vs.\ transmitted components.

\subsubsection{Panel D: Packet width $\sigma(t)$}

\paragraph{What is shown.}
The root-mean-square width of the packet,
\begin{equation}
  \sigma(t) \;=\; \sqrt{\langle x^2\rangle - \langle x\rangle^2},
  \quad
  \langle x^2\rangle = \dx\sum_i x_i^2\,|\psi_i|^2.
\end{equation}

\paragraph{Expected behaviour.}
\begin{itemize}
  \item \textbf{Free propagation}: A Gaussian wave packet in free space
        spreads as
        \begin{equation}
          \sigma(t) \;=\; \sigma_0\,\sqrt{1 + \left(\frac{t}{2\sigma_0^2}\right)^2},
          \label{eq:spreading}
        \end{equation}
        which for $t \gg 2\sigma_0^2 = 0.5$\,a.u.\ becomes linear:
        $\sigma(t) \approx t/(2\sigma_0)$.

  \item \textbf{After scattering}: The width grows faster than the free-space
        result.  The physical reason is that the wave function splits into
        two spatially separated components (reflected and transmitted), and
        the variance $\langle x^2\rangle - \langle x\rangle^2$ receives a
        large contribution from the \emph{distance between the two peaks},
        not just from the width of each individual peak.  This is a
        purely quantum interference effect.
\end{itemize}

\paragraph{What to look for.}
A clear change of slope at the time of scattering: the growth rate of
$\sigma(t)$ increases sharply when the packet splits.

\subsection{Animation (\texttt{evolution.mp4})}

The animation simultaneously shows $|\psi(x,t)|^2$ (solid blue) and
$\mathrm{Re}[\psi(x,t)]$ (dashed orange), which reveals the oscillatory
phase structure of the wave function.

\paragraph{Expected behaviour.}
\begin{itemize}
  \item The real part oscillates at spatial frequency $k_0$, visible as
        rapid oscillations inside the packet envelope.
  \item As the packet approaches the barrier, the oscillations in the
        barrier region are exponentially suppressed (evanescent wave).
  \item After the interaction, the reflected packet shows oscillations
        travelling in the $-x$ direction (decreasing phase with time),
        while the tiny transmitted packet oscillates in the $+x$ direction.
  \item Interference between the incoming and reflected components is
        visible as a standing-wave pattern to the left of the barrier
        during the scattering phase.
\end{itemize}

% ═══════════════════════════════════════════════════════════════════════════
\section{Numerical parameters: stability and accuracy}
% ═══════════════════════════════════════════════════════════════════════════

\subsection{Grid resolution requirements}

For the simulation to faithfully represent the wave packet, the grid
must resolve the shortest physical length scales:

\begin{itemize}
  \item \textbf{Spatial}: the de Broglie wavelength $\lambda = 2\pi/k_0$
        must be resolved with several grid points.  For $k_0 = 5$,
        $\lambda \approx 1.26$\,bohr, and $\dx = 0.01$\,bohr gives
        $\approx 126$ points per wavelength — more than sufficient.

  \item \textbf{Inside the barrier}: the evanescent decay length is
        $1/\kappa \approx 0.26$\,bohr for the default parameters.
        With $\dx = 0.01$, there are $\approx 26$ points per decay length,
        which is adequate.

  \item \textbf{Temporal}: the \CN{} scheme is unconditionally stable, so
        $\dt$ is constrained only by \emph{accuracy}.  A practical
        criterion is $\dt \ll 1/E_\text{max}$, where
        $E_\text{max} \approx k_\text{max}^2/2$ is the maximum energy in
        the spectrum.  For $k_\text{max} \approx k_0 + 3/(2\sigma)
        \approx 8$, $E_\text{max} \approx 32$\,a.u., giving
        $\dt \ll 0.03$\,a.u.  The default $\dt = 0.001$\,a.u.\ satisfies
        this easily.
\end{itemize}

\subsection{Courant number}

Although the \CN{} scheme is unconditionally stable, the effective
Courant number $\nu = k_0\dt/\dx = 5 \times 0.001 / 0.01 = 0.5$
gives an indication of the accuracy of the phase velocity.  Values
$\nu \lesssim 1$ ensure that the numerical group velocity closely
matches the analytical one.

% ═══════════════════════════════════════════════════════════════════════════
\section{Conclusions}
% ═══════════════════════════════════════════════════════════════════════════

The \CN{} scheme provides a robust, second-order, norm-conserving method
for propagating the TDSE.  Its implicit nature eliminates the time-step
restriction that affects explicit schemes, and the tridiagonal structure
of the resulting linear system makes each step $\mathcal{O}(N)$ via the
Thomas algorithm.

The simulation reproduces the key quantum features of barrier scattering:
wave-packet spreading, tunnelling through a classically forbidden region,
and the splitting of the packet into reflected and transmitted components.
The quantitative comparison of the numerically computed transmission
coefficient with the analytic plane-wave formula provides a direct
validation of both the physical model and the numerical implementation.

\clearpage

% ═══════════════════════════════════════════════════════════════════════════
\begin{thebibliography}{99}

\bibitem{Griffiths2018}
D.~J.~Griffiths and D.~F.~Schroeter,
\textit{Introduction to Quantum Mechanics}, 3rd ed.\
Cambridge University Press, 2018.

\bibitem{CrankNicolson1947}
J.~Crank and P.~Nicolson,
``A practical method for numerical evaluation of solutions of partial
differential equations of the heat-conduction type,''
\textit{Proc.\ Camb.\ Phil.\ Soc.}, vol.~43, pp.~50--67, 1947.

\bibitem{Thomas1949}
L.~H.~Thomas,
``Elliptic problems in linear difference equations over a network,''
Watson Sci.\ Comput.\ Lab.\ Report, Columbia University, 1949.

\bibitem{Tannor2007}
D.~J.~Tannor,
\textit{Introduction to Quantum Mechanics: A Time-Dependent Perspective},
University Science Books, 2007.

\bibitem{NumericalRecipes}
W.~H.~Press, S.~A.~Teukolsky, W.~T.~Vetterling, and B.~P.~Flannery,
\textit{Numerical Recipes: The Art of Scientific Computing}, 3rd ed.\
Cambridge University Press, 2007.

\bibitem{Sakurai2017}
J.~J.~Sakurai and J.~Napolitano,
\textit{Modern Quantum Mechanics}, 3rd ed.\
Cambridge University Press, 2021.

\bibitem{Cohen1977}
C.~Cohen-Tannoudji, B.~Diu, and F.~Lalo\"{e},
\textit{Quantum Mechanics}, vols.~I \& II,
Wiley, 1977.

\bibitem{Hairer1993}
E.~Hairer, S.~P.~N\o rsett, and G.~Wanner,
\textit{Solving Ordinary Differential Equations I: Nonstiff Problems},
2nd ed., Springer, 1993.

\bibitem{Leforestier1991}
C.~Leforestier et al.,
``A comparison of different propagation schemes for the time-dependent
Schr\"{o}dinger equation,''
\textit{J.\ Comput.\ Phys.}, vol.~94, pp.~59--80, 1991.

\bibitem{Wigner1955}
E.~P.~Wigner,
``Lower limit for the energy derivative of the scattering phase shift,''
\textit{Phys.\ Rev.}, vol.~98, pp.~145--147, 1955.

\end{thebibliography}

\clearpage

% ═══════════════════════════════════════════════════════════════════════════
\appendix
% ═══════════════════════════════════════════════════════════════════════════

% ───────────────────────────────────────────────────────────────────────────
\section{The quantum state space and Dirac formalism}
\label{app:hilbert}
% ───────────────────────────────────────────────────────────────────────────

\subsection{Hilbert space structure}

The rigorous mathematical setting for quantum mechanics is a
\emph{complex separable Hilbert space} $\mathcal{H}$.  For a
one-dimensional spinless particle, the natural choice is
$\mathcal{H} = L^2(\mathbb{R})$: the space of square-integrable
complex functions on $\mathbb{R}$, equipped with the inner product
\begin{equation}
  \langle f, g\rangle \;=\; \int_{-\infty}^{+\infty} f^*(x)\,g(x)\,dx.
\end{equation}
A physical state corresponds to a \emph{unit vector} (ray) in
$\mathcal{H}$, i.e.\ $\|\psi\| = 1$.

In Dirac notation, states are \emph{kets} $\ket{\psi} \in \mathcal{H}$
and their duals are \emph{bras} $\bra{\psi} \in \mathcal{H}^*$.
The inner product is $\braket{\phi}{\psi}$, and the norm is
$\|\psi\|^2 = \braket{\psi}{\psi}$.

\subsection{Observables as self-adjoint operators}

Every physical observable $A$ corresponds to a \emph{self-adjoint}
(Hermitian) operator $\hat{A}: \mathcal{H} \to \mathcal{H}$, defined
by $\hat{A}^\dagger = \hat{A}$.  Self-adjointness guarantees:
\begin{itemize}
  \item All eigenvalues are real (measurable quantities must be real).
  \item Eigenvectors corresponding to distinct eigenvalues are orthogonal.
  \item The eigenvectors form a complete orthonormal basis of $\mathcal{H}$
        (spectral theorem).
\end{itemize}

For continuous spectra (position, momentum), the eigenstates are
distributions rather than true $L^2$ functions.  The position
eigenstates $\ket{x'}$ satisfy
\begin{equation}
  \hat{x}\ket{x'} = x'\ket{x'}, \qquad
  \braket{x}{x'} = \delta(x-x'), \qquad
  \int_{-\infty}^{+\infty} \ket{x'}\bra{x'}\,dx' = \hat{I},
\end{equation}
and the wave function is the projection
$\psi(x,t) = \braket{x}{\psi(t)}$.

\subsection{The canonical commutation relation}

The position and momentum operators satisfy
\begin{equation}
  [\hat{x},\, \hat{p}] \;=\; \ii\hbar\,\hat{I}.
  \label{eq:ccr}
\end{equation}
In the position representation, $\hat{p} = -\ii\hbar\,\partial_x$.
The commutation relation \eqref{eq:ccr} is the algebraic root of the
Heisenberg uncertainty principle: for any state $\ket{\psi}$,
\begin{equation}
  \Delta x\,\Delta p \;\geq\; \frac{\hbar}{2},
\end{equation}
where $(\Delta A)^2 = \langle \hat{A}^2\rangle - \langle\hat{A}\rangle^2$.
The Gaussian wave packet \eqref{eq:gauss} saturates this bound, making
it the \emph{coherent state} of the Weyl--Heisenberg algebra.

\subsection{Time evolution operator and unitarity}

The formal solution of the TDSE is
\begin{equation}
  \ket{\psi(t)} \;=\; \hat{U}(t,t_0)\,\ket{\psi(t_0)},
  \qquad
  \hat{U}(t,t_0) = \exp\!\left[-\frac{\ii}{\hbar}\hat{H}(t-t_0)\right],
\end{equation}
where the last expression holds for time-independent $\hat{H}$.
The operator $\hat{U}$ is \emph{unitary}: $\hat{U}^\dagger \hat{U} = \hat{I}$,
which is the operator statement of probability conservation.
The \CN{} propagator $A^{-1}B$ is a finite-dimensional unitary
approximation to $\hat{U}(\dt)$.

% ───────────────────────────────────────────────────────────────────────────
\section{Transfer matrix method and exact tunnelling formulae}
\label{app:transfer}
% ───────────────────────────────────────────────────────────────────────────

\subsection{Piecewise solution of the stationary Schr\"{o}dinger equation}

For a static potential the TDSE separates via $\psi(x,t) =
\phi(x)\,\ee^{-\ii Et/\hbar}$, yielding the
\emph{time-independent Schr\"{o}dinger equation} (TISE):
\begin{equation}
  -\frac{1}{2}\phi''(x) + V(x)\phi(x) \;=\; E\,\phi(x).
\end{equation}

For the rectangular barrier \eqref{eq:barrier}, the TISE has piecewise
constant coefficients.  In each region the general solution is:
\begin{align}
  \text{Region I } (x < x_\text{bar}):&\quad
    \phi(x) = A\,\ee^{\ii kx} + B\,\ee^{-\ii kx}, \quad k = \sqrt{2E}, \\
  \text{Region II } (x_\text{bar} < x < x_\text{bar}+w):&\quad
    \phi(x) = C\,\ee^{\kappa x} + D\,\ee^{-\kappa x}, \quad
    \kappa = \sqrt{2(V_0-E)},\quad (E<V_0), \\
  \text{Region III } (x > x_\text{bar}+w):&\quad
    \phi(x) = F\,\ee^{\ii kx},
\end{align}
where $F\ee^{\ii kx}$ is the transmitted wave (no left-moving component
in region III for a right-incident wave).

\subsection{The transfer matrix}

At each interface, continuity of $\phi$ and $\phi'$ provides two
matching conditions.  These can be compactly encoded in a
$2\times 2$ \emph{transfer matrix} $M$ such that
\begin{equation}
  \begin{pmatrix} A \\ B \end{pmatrix}
  \;=\; M \begin{pmatrix} F \\ 0 \end{pmatrix}.
\end{equation}
For the rectangular barrier, $M = M_{\text{I}\to\text{II}}\cdot
M_{\text{II}\to\text{III}}$, where each interface matrix is
\begin{equation}
  M_{j\to k} \;=\; \frac{1}{2q_k}
  \begin{pmatrix}
    q_k + q_j & q_k - q_j \\
    q_k - q_j & q_k + q_j
  \end{pmatrix},
\end{equation}
with $q_j$ being the (possibly imaginary) wave vector in region $j$.
The transmission amplitude is $t = 1/M_{11}$ and
\begin{equation}
  T = |t|^2 = \frac{1}{|M_{11}|^2},
\end{equation}
which, after explicit computation, reproduces \eqref{eq:T_tunnelling}.
The transfer matrix formalism generalises straightforwardly to
arbitrary piecewise-constant potentials.

\subsection{WKB approximation}

For a \emph{smooth} potential barrier, the exact transfer matrix is
not available in closed form.  The Wentzel--Kramers--Brillouin (WKB)
approximation replaces the exact evanescent solution inside the barrier
with a locally-plane-wave ansatz:
\begin{equation}
  \phi(x) \;\approx\;
  \frac{C}{\sqrt{\kappa(x)}}\,\exp\!\left[-\int_{x_1}^{x}\kappa(x')\,dx'\right],
  \quad \kappa(x) = \sqrt{2(V(x)-E)},
\end{equation}
where $x_1$, $x_2$ are the classical turning points ($V(x_{1,2}) = E$).
The WKB transmission coefficient is
\begin{equation}
  T_\text{WKB} \;\approx\; \exp\!\left[-2\int_{x_1}^{x_2}\kappa(x)\,dx\right].
  \label{eq:wkb}
\end{equation}
For the rectangular barrier, $\kappa(x) = \kappa = \text{const}$ and
\eqref{eq:wkb} gives $T \approx \ee^{-2\kappa w}$, consistent with the
large-$\kappa w$ limit of \eqref{eq:T_tunnelling}.

% ───────────────────────────────────────────────────────────────────────────
\section{Dispersion relation and free-packet dynamics}
\label{app:dispersion}
% ───────────────────────────────────────────────────────────────────────────

\subsection{Dispersion relation}

For a free particle ($V=0$), substituting a plane wave
$\psi = \ee^{\ii(kx - \omega t)}$ into the TDSE gives the
\emph{dispersion relation}
\begin{equation}
  \omega(k) \;=\; \frac{k^2}{2} \quad \text{(atomic units)}.
\end{equation}
This is a \emph{quadratic} dispersion, unlike the linear dispersion of
photons.  Two velocities characterise wave-packet propagation:
\begin{align}
  v_\text{phase} &= \frac{\omega}{k} = \frac{k}{2}, \\
  v_\text{group} &= \frac{d\omega}{dk} = k.
\end{align}
The group velocity $v_g = k_0$ is the velocity of the packet's centre
of mass and equals the classical velocity $p/m = k_0$ in atomic units.

\subsection{Exact free-propagation of the Gaussian packet}

A Gaussian wave packet in free space can be propagated exactly using
the Fourier convolution theorem.  Starting from \eqref{eq:gauss}, its
Fourier transform is
\begin{equation}
  \tilde{\phi}(k,0) \;=\;
  (2\pi\sigma^2)^{1/4}\,\sigma\,\ee^{-(k-k_0)^2\sigma^2}\,\ee^{-\ii k x_0}.
\end{equation}
At time $t$, each mode acquires a phase $\ee^{-\ii\omega(k)t}$:
\begin{equation}
  \tilde{\phi}(k,t) \;=\; \tilde{\phi}(k,0)\,\ee^{-\ii k^2 t/2}.
\end{equation}
Transforming back gives
\begin{equation}
  \psi(x,t) \;=\;
  \left(\frac{1}{2\pi\sigma_t^2}\right)^{1/4}
  \exp\!\left[-\frac{(x - x_0 - k_0 t)^2}{4\sigma_t^2}\right]
  \exp\!\left[\ii k_0\!\left(x - \frac{k_0 t}{2}\right)\right]
  \exp\!\left[\ii\phi_t\right],
  \label{eq:free_packet}
\end{equation}
where the \emph{complex width} is
\begin{equation}
  \sigma_t \;=\; \sigma\sqrt{1 + \ii t/(2\sigma^2)},
\end{equation}
and $\phi_t = -\frac{1}{2}\arg(\sigma_t/\sigma)$ is a global phase.
The probability density is still Gaussian:
\begin{equation}
  |\psi(x,t)|^2 \;=\;
  \frac{1}{\sqrt{2\pi}|\sigma_t|}
  \exp\!\left[-\frac{(x-x_0-k_0 t)^2}{2|\sigma_t|^2}\right],
\end{equation}
with real width $|\sigma_t| = \sigma\sqrt{1 + t^2/(4\sigma^4)}$,
which is exactly equation \eqref{eq:spreading}.

The spreading rate $d|\sigma_t|/dt \to 1/(2\sigma)$ for large $t$
is inversely proportional to the initial width: a narrow packet
(small $\sigma$, large momentum spread) spreads faster.

% ───────────────────────────────────────────────────────────────────────────
\section{Operator splitting and the Cayley form}
\label{app:splitting}
% ───────────────────────────────────────────────────────────────────────────

\subsection{The exact propagator as a matrix exponential}

The formal solution of the TDSE over one time step is
\begin{equation}
  \boldsymbol{\psi}^{n+1} \;=\;
  \ee^{-\ii\dt H_\text{disc}}\,\boldsymbol{\psi}^n.
\end{equation}
Direct computation of the matrix exponential costs $\mathcal{O}(N^3)$
— prohibitive for large $N$.  Several efficient approximations exist.

\subsection{Padé approximants and the Cayley form}

The \CN{} scheme is equivalent to the \emph{first-order Padé
approximant} of the matrix exponential:
\begin{equation}
  \ee^{-\ii\dt H} \;\approx\;
  \frac{I - \ii\dt H/2}{I + \ii\dt H/2}
  \;=\; A^{-1}B.
  \label{eq:cayley}
\end{equation}
This is also known as the \emph{Cayley form} of the propagator.  It is
\emph{exactly unitary} regardless of $\dt$, because the numerator and
denominator are complex conjugates of each other (for Hermitian $H$).

Higher-order Padé approximants $[m/m]$ exist and achieve
$\mathcal{O}(\dt^{2m})$ accuracy while remaining unitary, but they
require solving $m$ tridiagonal systems per step and are rarely used
in practice.

\subsection{Strang splitting}

For Hamiltonians of the form $H = T + V$ where $T$ is the kinetic
operator and $V$ is diagonal in position space, the
\emph{split-operator} (Strang splitting) method approximates:
\begin{equation}
  \ee^{-\ii\dt(T+V)} \;\approx\;
  \ee^{-\ii\dt V/2}\,\ee^{-\ii\dt T}\,\ee^{-\ii\dt V/2}
  + \mathcal{O}(\dt^3).
\end{equation}
Each factor can be applied efficiently: $\ee^{-\ii\dt V/2}$ is a
pointwise multiplication in position space (diagonal), while
$\ee^{-\ii\dt T}$ is a multiplication in momentum (Fourier) space.
Using the Fast Fourier Transform (FFT), each step costs
$\mathcal{O}(N\log N)$, compared to $\mathcal{O}(N)$ for the \CN{}
Thomas algorithm.  However, the split-operator method is second-order
accurate and exactly unitary only in the limit of exact arithmetic, and
requires global FFTs which are harder to parallelise.

% ───────────────────────────────────────────────────────────────────────────
\section{Scattering theory: S-matrix and Wigner delay time}
\label{app:scattering}
% ───────────────────────────────────────────────────────────────────────────

\subsection{The scattering matrix}

For a potential that vanishes at $x \to \pm\infty$, scattering theory
describes the asymptotic relation between incoming and outgoing plane
waves.  In 1D, the \emph{S-matrix} is a $2\times 2$ unitary matrix:
\begin{equation}
  S \;=\;
  \begin{pmatrix}
    t   & r' \\
    r   & t'
  \end{pmatrix},
\end{equation}
where $t, t'$ are transmission amplitudes (right-incident and
left-incident) and $r, r'$ are reflection amplitudes.  Unitarity
$S^\dagger S = I$ implies $|t|^2 + |r|^2 = 1$, i.e.\ $T + R = 1$.

For a symmetric potential (even $V(x)$), $t = t'$ and $|r| = |r'|$.
The rectangular barrier is symmetric about its centre, so both
transmission amplitudes are equal.

\subsection{Phase shifts and time delays}

The transmission amplitude can be written as
$t = \sqrt{T}\,\ee^{\ii\varphi_t(k)}$, where $\varphi_t$ is the
transmission phase.  The \emph{Wigner--Smith delay time} is defined as
\begin{equation}
  \tau(k) \;=\; \frac{d\varphi_t}{d\omega}
  \;=\; \frac{1}{v_g}\frac{d\varphi_t}{dk},
  \label{eq:wigner}
\end{equation}
and represents the extra time spent by the transmitted wave packet
near the barrier compared to free propagation \cite{Wigner1955}.
For the rectangular barrier in the tunnelling regime,
$\tau > 0$ (the packet is delayed), and as $V_0 \to \infty$ the
delay saturates to $\tau_\text{max} = w/v_g$ — the so-called
\emph{Hartman effect}, where the delay becomes independent of
barrier width for opaque barriers.

In the simulation, the delay manifests as a reduction in the slope
of $\langle x\rangle(t)$ during the scattering event (Panel C of
\texttt{analysis.png}).

\subsection{The optical theorem in 1D}

The optical theorem in 1D relates the forward scattering amplitude to
the total cross-section.  In one dimension, ``cross-section'' is
replaced by the \emph{total reflectance}, and the theorem states:
\begin{equation}
  \mathrm{Im}(r) + \frac{k}{2}(T-1) = 0,
\end{equation}
which is a direct consequence of $S$-matrix unitarity and
provides an additional consistency check for numerical scattering
calculations.

% ───────────────────────────────────────────────────────────────────────────
\section{Von Neumann stability analysis}
\label{app:stability}
% ───────────────────────────────────────────────────────────────────────────

\subsection{Von Neumann analysis for the TDSE}

The stability of finite-difference schemes for linear PDEs can be
analysed rigorously via the von Neumann (Fourier) method.  Consider a
uniform grid and a single Fourier mode $\psi_j^n = g^n\,\ee^{\ii j\xi\dx}$,
where $\xi$ is the wave vector and $g$ is the amplification factor.
Stability requires $|g| \leq 1$ for all $\xi$.

\paragraph{Forward Euler (explicit):}
Substituting into $\psi_j^{n+1} = \psi_j^n - \ii\dt H_\text{disc}\psi_j^n$
gives
\begin{equation}
  g \;=\; 1 - \ii\dt\,\lambda(\xi),
  \quad \lambda(\xi) = \frac{2(1-\cos(\xi\dx))}{\dx^2} + V \in \mathbb{R}.
\end{equation}
Then $|g|^2 = 1 + \dt^2\lambda^2 > 1$ for $\lambda \neq 0$: the scheme
is \emph{always unstable} for the Schrödinger equation.

\paragraph{Crank--Nicolson:}
The amplification factor is
\begin{equation}
  g \;=\; \frac{1 - \ii\dt\lambda/2}{1 + \ii\dt\lambda/2}.
\end{equation}
Since $\lambda \in \mathbb{R}$, we have $|g|^2 = 1$ identically for
all $\xi$ and all $\dt > 0$: the scheme is \emph{unconditionally
stable} and, moreover, \emph{exactly norm-preserving} at the level of
the amplification factor.

\paragraph{Implicit mid-point rule (Crank--Nicolson is a special case):}
More generally, the $\theta$-method
\begin{equation}
  \ii\frac{\psi^{n+1}-\psi^n}{\dt} \;=\;
  (1-\theta)H\psi^n + \theta H\psi^{n+1}
\end{equation}
is stable for $\theta \geq 1/2$.  The Crank--Nicolson scheme
corresponds to $\theta = 1/2$ and is the only choice in this family
that is both stable and second-order accurate.

\subsection{Numerical dispersion}

Although the \CN{} scheme is unconditionally stable, it introduces
a \emph{numerical dispersion error}: the discrete phase velocity
$\tilde{\omega}(\xi)$ differs from the exact $\omega(\xi) = \xi^2/2$.
From the amplification factor,
\begin{equation}
  \tilde{\omega}(\xi) \;=\;
  \frac{2}{\dt}\arctan\!\left(\frac{\dt\,\lambda(\xi)}{2}\right),
  \quad
  \lambda(\xi) = \frac{2(1-\cos\xi\dx)}{\dx^2}.
\end{equation}
Expanding for small $\xi\dx$ and $\dt$:
\begin{equation}
  \tilde{\omega}(\xi) \;=\; \frac{\xi^2}{2}
  \left[1 - \frac{(\xi\dx)^2}{12} - \frac{\dt^2\xi^4}{24} + \ldots\right],
\end{equation}
showing $\mathcal{O}(\dx^2, \dt^2)$ dispersion errors, consistent with
the second-order accuracy of the scheme.  High-$k$ components of the
wave packet travel at a slightly wrong speed — this is the physical
origin of the \emph{numerical spreading} that adds to the physical
quantum spreading.

% ───────────────────────────────────────────────────────────────────────────
\section{Ehrenfest's theorem and quantum-classical correspondence}
\label{app:ehrenfest}
% ───────────────────────────────────────────────────────────────────────────

\subsection{Statement and derivation}

\emph{Ehrenfest's theorem} establishes the connection between quantum
expectation values and classical equations of motion.  For any
observable $\hat{A}$, the Heisenberg equation of motion is:
\begin{equation}
  \frac{d}{dt}\langle\hat{A}\rangle \;=\;
  \frac{1}{\ii\hbar}\langle[\hat{A},\hat{H}]\rangle
  + \left\langle\pdv{\hat{A}}{t}\right\rangle.
\end{equation}
Applying this to $\hat{x}$ and $\hat{p}$ with $\hat{H} = \hat{p}^2/2m + V(\hat{x})$:
\begin{align}
  \frac{d\langle\hat{x}\rangle}{dt}
  &= \frac{\langle\hat{p}\rangle}{m}, \label{eq:ehr1} \\
  \frac{d\langle\hat{p}\rangle}{dt}
  &= -\left\langle\frac{\partial V}{\partial x}\right\rangle. \label{eq:ehr2}
\end{align}
These are the quantum analogues of Hamilton's equations.  Note
the crucial difference from the classical result: \eqref{eq:ehr2}
involves $\langle V'(\hat{x})\rangle$, not $V'(\langle\hat{x}\rangle)$.
The two coincide only if $V$ is at most quadratic, or if the wave
packet is so narrow that $\psi(x,t) \approx \delta(x - x_\text{cl}(t))$.

\subsection{Implications for the simulation}

In the early propagation phase, the wave packet is narrow compared to
the scale on which $V$ varies (the barrier), so Ehrenfest's theorem
predicts nearly classical motion: $\langle x\rangle(t) \approx
x_0 + k_0 t$.  This is exactly the linear regime visible in Panel C
before the scattering event.

When the packet interacts with the barrier, $\langle V'\rangle \neq
V'(\langle x\rangle)$, the classical approximation breaks down, and
quantum effects (tunnelling, interference) dominate.  After
the scattering, the wave function consists of two well-separated
packets, and $\langle x\rangle$ becomes a weighted average that has
no classical counterpart.

\subsection{Energy conservation}

From Ehrenfest's theorem, the mean energy $\langle\hat{H}\rangle$ is
conserved:
\begin{equation}
  \frac{d}{dt}\langle\hat{H}\rangle \;=\;
  \frac{1}{\ii\hbar}\langle[\hat{H},\hat{H}]\rangle = 0.
\end{equation}
In the numerical simulation, energy conservation is not enforced
explicitly by the \CN{} scheme (only norm is exactly conserved).
Nevertheless, the second-order accuracy of the scheme ensures that
the energy drift is $\mathcal{O}(\dt^2)$ per step.

% ───────────────────────────────────────────────────────────────────────────
\section{Absorbing boundary conditions and perfectly matched layers}
\label{app:abc}
% ───────────────────────────────────────────────────────────────────────────

\subsection{The problem with Dirichlet boundaries}

Dirichlet boundary conditions $\psi(x_\text{min}) = \psi(x_\text{max}) = 0$
mimic infinitely hard walls.  They are simple to implement but introduce
\emph{spurious reflections} as soon as the wave packet reaches the boundary.
The domain must therefore be large enough that no part of $\psi$ reaches
the walls during the simulation — a requirement that can make the problem
expensive for long times or fast packets.

\subsection{Absorbing boundary conditions}

A more efficient alternative is to modify the Hamiltonian near the
boundaries by adding a negative imaginary potential (absorption):
\begin{equation}
  V(x) \;\to\; V(x) - \ii\,\Gamma(x),
  \quad \Gamma(x) \geq 0,
\end{equation}
where $\Gamma(x)$ is large only near the domain edges.  This makes
the Hamiltonian \emph{non-Hermitian} in the absorbing region, causing
the norm to decay as probability is absorbed.  The reflected numerical
wave is exponentially suppressed.

A standard choice is the cosine ramp:
\begin{equation}
  \Gamma(x) \;=\;
  \Gamma_0\,\sin^2\!\left(\frac{\pi}{2}
  \cdot\frac{|x|-x_\text{abs}}{L_\text{abs}}\right)
  \quad \text{for } |x| > x_\text{abs},
\end{equation}
where $L_\text{abs}$ and $\Gamma_0$ are the absorption layer width and
strength.

\subsection{Perfectly matched layers}

The gold standard for absorbing boundaries is the \emph{perfectly
matched layer} (PML), originally developed for Maxwell's equations
\cite{NumericalRecipes} and later adapted to the TDSE.  A PML is
constructed by analytically continuing the spatial coordinate into the
complex plane:
\begin{equation}
  x \;\to\; \tilde{x}(x) \;=\; x + \ii\int_0^x \sigma(x')\,dx',
\end{equation}
where $\sigma(x') \geq 0$ is the PML absorption profile.  This
transformation preserves the form of the Schr\"{o}dinger equation in
the PML region while making all outgoing waves evanescent — and
crucially, \emph{without creating any reflection at the PML interface},
regardless of the angle or frequency of the incoming wave.

For the 1D TDSE, the PML modification is equivalent to introducing a
complex coordinate stretching factor $s(x)$:
\begin{equation}
  \Htop_\text{PML} \;=\;
  -\frac{1}{2s(x)}\frac{\partial}{\partial x}\frac{1}{s(x)}\frac{\partial}{\partial x}
  + V(x),
  \quad s(x) = 1 + \ii\sigma(x)/\omega.
\end{equation}
The present simulation uses simple Dirichlet BCs, which are sufficient
for the default parameters.  Implementing absorbing BCs would allow
simulations on much smaller domains or over much longer times.

\end{document}
